\documentclass[11pt]{amsart}
\usepackage[T1]{fontenc}
\usepackage[utf8]{inputenc}
\usepackage[spanish,es-noquoting]{babel}
\usepackage{lmodern}
\usepackage{amsmath,amssymb,mathtools}
\usepackage{booktabs,array,longtable}
\usepackage{tikz}
\usetikzlibrary{arrows.meta,positioning,calc}
\usepackage{url}
\usepackage{enumitem}
\usepackage{float}
\usepackage[hidelinks]{hyperref}
\hypersetup{
  pdftitle={Comprobación y ampliación de la hipótesis de Riemann},
  pdfauthor={Germán García Andrés},
  pdfsubject={Función zeta y xi de Riemann; involuciones; familias F2/F4; grafos espectrales; homología persistente; criterio de celdas unitarias; operador monopolar y correspondencia estructural},
  pdfkeywords={Hipótesis de Riemann, función xi, ceros no triviales, simplicidad, involución, grafos, Laplaciano, homología persistente, Vietoris-Rips, operador de simetría, operador monopolar, matrices aleatorias, espectros cuánticos},
  pdfcreator={LaTeX / AMS article class},
  pdfproducer={pdfTeX}
}

\newtheorem{theorem}{Teorema}[section]
\newtheorem{proposition}[theorem]{Proposición}
\newtheorem{lemma}[theorem]{Lema}
\newtheorem{corollary}[theorem]{Corolario}
\theoremstyle{definition}
\newtheorem{definition}[theorem]{Definición}
\newtheorem{algorithmicdef}[theorem]{Procedimiento}
\theoremstyle{remark}
\newtheorem{remark}[theorem]{Observación}
\newtheorem{scope}[theorem]{Condición de alcance}

\title[Comprobación y ampliación de la hipótesis de Riemann]{Comprobación y ampliación de la hipótesis de Riemann}
\author{Germán García Andrés}
\thanks{Investigador independiente}
\address{Investigador independiente}
\date{26 de agosto de 2026}
\subjclass[2020]{11M26, 05C50, 55N31}
\keywords{Hipótesis de Riemann, función xi, simetría, involución, órbitas, grafos, Laplaciano, operador de simetría, equivarianza, homología persistente, superzeta, familias generacionales, operador monopolar, breathing mode, correspondencia estructural, matrices aleatorias, espectros cuánticos}

\begin{document}
\begin{abstract}
Se presenta un marco analítico--topológico para estudiar los ceros no triviales de la función zeta de Riemann mediante involuciones, órbitas, operadores inducidos y grafos complementarios. La involución crítica $R(s)=1-\overline{s}$ proporciona la caracterización exacta $R(\rho)=\rho\iff\Re(\rho)=1/2$, y su representación $P_R$ sobre el grafo satisface $P_R^2=I$ y, bajo preservación de aristas, $[L,P_R]=0$. Se formaliza la clasificación orbital $F_2/F_4$ y un criterio reforzado de celdas unitarias: la existencia exhaustiva de regiones $R$-invariantes $B_n$ con $N(B_n)=1$ es equivalente a RH junto con la simplicidad de todos los ceros no triviales. El método se complementa con controles nulo, neutro y no nulo, una filtración de Vietoris--Rips y una implementación reproducible para la estructura $H_0$ de las ordenadas desplegadas. Como extensión operatoria se desarrolla el modo monopolar radial $L=0$ y su Hamiltoniano efectivo, y se formula una correspondencia estructural entre las líneas analítico--gráfica y operatorio--espectral basada en involución, sector fijo, conmutación y espectro. La relación se sitúa en el contexto de las conexiones Montgomery--Dyson entre ceros de Riemann, matrices aleatorias y espectros cuánticos.
\end{abstract}
\maketitle

\section{Objeto y alcance}
La Hipótesis de Riemann (RH) afirma que todo cero no trivial $\rho$ de $\zeta(s)$ satisface
\begin{equation}\label{eq:RH}
   \Re(\rho)=\frac12.
\end{equation}
La formulación oficial del Clay Mathematics Institute es precisamente la afirmación universal \eqref{eq:RH}. El objetivo de este manuscrito es exponer, con detalle reproducible, el método desarrollado para: (i) expresar el cierre de la simetría en términos de involuciones y órbitas; (ii) clasificar las órbitas en familias $F_2$ y $F_4$; (iii) construir grafos nulo, neutro y no nulo; (iv) identificar agrupamientos con Laplaciano y persistencia $H_0$; (v) representar la simetría por un operador $P_R$ y estudiar su conmutación con el Laplaciano; (vi) organizar una jerarquía generacional; y (vii) formular un criterio reforzado que, al certificarse globalmente, produce RH junto con simplicidad.

El cuerpo matemático es complejo-analítico, topológico y espectral. La sección monopolar desarrolla, en paralelo, una extensión operatoria física basada en el modo radial $L=0$ y conserva explícitamente la separación entre el operador nuclear y el operador de simetría del grafo.

\section{Mapa conceptual y resultados principales}\label{sec:mainresults}
La arquitectura completa puede leerse de izquierda a derecha como
\begin{equation}\label{eq:roadmap}
\boxed{\begin{aligned}
\xi &\longrightarrow (C,F,R)\longrightarrow F_2/F_4\longrightarrow P_R,\\
P_R &\longrightarrow [L,P_R]\longrightarrow B_n\longrightarrow N(B_n)=1\\
&\longrightarrow \mathrm{RH}+\text{simplicidad}.
\end{aligned}}
\end{equation}
Cada flecha tiene un papel diferente: las primeras describen la simetría analítica y su representación gráfica; las intermedias introducen el lenguaje espectral; y la última aplica el principio del argumento a regiones invariantes.

\subsection{Resultados enunciados al inicio}
\begin{enumerate}[label=\textbf{R\arabic*.},leftmargin=*]
\item \textbf{Caracterización crítica exacta.} Para todo cero no trivial $\rho$,
\[
\Delta_R(\rho)=0\iff R(\rho)=\rho\iff F(\rho)=C(\rho)\iff \Re(\rho)=\frac12.
\]
La demostración aparece en la Proposición~\ref{prop:residual}.
\item \textbf{Representación operatoria.} La asignación de ceros a vértices es equivariante,
\[
Q\circ R=P_R\circ Q,
\]
y, si $P_R$ preserva las aristas del grafo, entonces $P_R^2=I$ y $[L,P_R]=0$; véanse las Proposiciones~\ref{prop:equivariance} y~\ref{prop:lapcomm}.
\item \textbf{Criterio global.} RH equivale a que todo cero no trivial pertenezca al sector fijo de $R$, o equivalentemente a $\Delta_R(\rho)=0$ para todo $\rho\in Z^*$.
\item \textbf{Ampliación a simplicidad.} El Teorema~\ref{thm:unitcells} establece la equivalencia entre RH más simplicidad de todos los ceros y la existencia, para cada cero, de una región $R$-invariante con índice unitario $N(B_n)=1$.
\item \textbf{Cierre por saturación de conjuntos.} Para cada altura finita $T$, una inclusión $S(T)\subseteq Z_2(T)\subseteq Z^*(T)$ junto con la igualdad de conteos $|S(T)|=N(T)=\sum_{\rho\in Z^*(T)}m(\rho)$ fuerza $S(T)=Z_2(T)=Z^*(T)$; por tanto $Z_4(T)=\varnothing$. Una sucesión exhaustiva $T_k\to\infty$ transmite el cierre a todo el strip crítico.
\end{enumerate}

\subsection{Estatus de las afirmaciones}
Para facilitar la lectura se separan cuatro niveles de contenido:
\begin{center}
\begin{tabular}{p{0.25\linewidth}p{0.66\linewidth}}
\toprule
\textbf{Nivel} & \textbf{Contenido}\\
\midrule
Resultado exacto & Identidades de simetría, caracterización de $\operatorname{Fix}(R)$, clasificación $F_2/F_4$, equivarianza y conmutación bajo automorfismo.\\
Teorema derivado & Criterio reforzado de celdas unitarias y su equivalencia con RH más simplicidad.\\
Resultado computacional & Grafos nulo--neutro--no nulo, espectros laplacianos y persistencia $H_0$ en bloques finitos de ordenadas.\\
Extensión operatoria & Operador monopolar nuclear $L=0$, Hamiltoniano colectivo y comparación de arquitectura con el operador gráfico.\\
\bottomrule
\end{tabular}
\end{center}

\subsection{Tabla mínima de símbolos}
\begin{center}
\begin{tabular}{p{0.18\linewidth}p{0.72\linewidth}}
\toprule
\textbf{Símbolo} & \textbf{Significado}\\
\midrule
$C(s)$ & Conjugación compleja $\overline{s}$.\\
$F(s)$ & Reflexión funcional $1-s$.\\
$R(s)$ & Involución crítica $1-\overline{s}$.\\
$\Delta_R$ & Residuo de cierre $F(\rho)-C(\rho)=1-2\Re(\rho)$.\\
$F_2,F_4$ & Órbitas geométricas crítica y genérica, respectivamente.\\
$Q$ & Asignación de un cero a su vértice; en la capa generacional también se usa un cociente escalar, indicado explícitamente en contexto.\\
$P_R$ & Operador de permutación inducido por $R$ sobre vértices.\\
$A,D,L$ & Matrices de adyacencia, grados y Laplaciano $L=D-A$.\\
$\beta_0$ & Número de componentes conexas, $\dim\ker L$.\\
$B_n$ & Región de Jordan $R$-invariante asociada al cero $\rho_n$.\\
$N(B)$ & Número de ceros de $\xi$ en $B$, contados con multiplicidad.\\
$\mathcal J$ & Involución radial $x\mapsto -x$ en la extensión monopolar.\\
$\widehat H_M$ & Hamiltoniano colectivo cuantizado del modo monopolar ideal.\\
\bottomrule
\end{tabular}
\end{center}

\section{Preliminares analíticos establecidos}
\subsection{Zeta, ceros triviales y función completada}
La función $\zeta(s)$, inicialmente definida para $\Re(s)>1$ por
\begin{equation}
  \zeta(s)=\sum_{n=1}^{\infty}n^{-s},
\end{equation}
se prolonga meromórficamente al plano complejo con un polo simple en $s=1$. La ecuación funcional implica los ceros triviales
\begin{equation}
  -2,-4,-6,\ldots.
\end{equation}
Los demás ceros se encuentran en la franja crítica $0<\Re(s)<1$ y se denominan no triviales.

Es útil trabajar con la función entera de Riemann
\begin{equation}\label{eq:xi}
 \xi(s)=\frac12 s(s-1)\pi^{-s/2}\Gamma\!\left(\frac{s}{2}\right)\zeta(s),
\end{equation}
que satisface
\begin{equation}\label{eq:functionalxi}
   \xi(s)=\xi(1-s)
\end{equation}
y cuyos ceros son exactamente los ceros no triviales de $\zeta$ contando multiplicidades.

\subsection{Conjugación y reflexión}
Definimos
\begin{equation}\label{eq:CFR}
  C(s)=\overline{s},\qquad F(s)=1-s,\qquad R(s)=F(C(s))=1-\overline{s}.
\end{equation}
Las tres transformaciones son involutivas:
\begin{equation}
 C^2=F^2=R^2=I.
\end{equation}
Además, la realidad de los coeficientes y \eqref{eq:functionalxi} hacen que estas transformaciones preserven el conjunto de ceros de $\xi$.

\section{Familias de simetría: $F_2$ y $F_4$}
Sea
\begin{equation}
 \rho=\sigma+i\gamma
\end{equation}
un cero no trivial. Consideremos la órbita bajo las transformaciones generadas por $C$ y $F$:
\begin{equation}\label{eq:orbit}
 \mathcal O(\rho)=\{\rho,\overline\rho,1-\rho,1-\overline\rho\}.
\end{equation}

\begin{definition}[Familia $F_4$]
Si $\sigma\neq1/2$ y $\gamma\neq0$, los cuatro elementos de \eqref{eq:orbit} son genéricamente distintos. Denotamos esta órbita por $F_4$.
\end{definition}

\begin{definition}[Familia $F_2$]
Si $\sigma=1/2$, entonces
\begin{equation}\label{eq:F2collapse}
  1-\rho=\overline\rho,\qquad 1-\overline\rho=\rho,
\end{equation}
y la órbita se reduce a dos posiciones geométricas:
\begin{equation}
 F_2(\gamma)=\left\{\frac12+i\gamma,\frac12-i\gamma\right\}.
\end{equation}
\end{definition}

\begin{proposition}[Residuo de cierre reversible]\label{prop:residual}
Defina
\begin{equation}\label{eq:delta}
 \Delta_R(\rho)=F(\rho)-C(\rho).
\end{equation}
Entonces
\begin{equation}\label{eq:deltaeval}
 \Delta_R(\rho)=1-2\Re(\rho),
\end{equation}
y por tanto
\begin{equation}\label{eq:equiv}
 \Delta_R(\rho)=0
 \iff F(\rho)=C(\rho)
 \iff R(\rho)=\rho
 \iff \Re(\rho)=\frac12.
\end{equation}
\end{proposition}

\begin{proof}
Para $\rho=\sigma+i\gamma$,
\[
F(\rho)-C(\rho)=1-\sigma-i\gamma-(\sigma-i\gamma)=1-2\sigma.
\]
Así, $\Delta_R=0$ si y sólo si $\sigma=1/2$. Además,
$R(\rho)=\rho$ equivale a $1-\overline\rho=\rho$, que produce la misma igualdad real $1-\sigma=\sigma$.
\end{proof}

\begin{corollary}[Conjunto fijo]
\begin{equation}
 \operatorname{Fix}(R)=\{s\in\mathbb C:\Re(s)=1/2\}.
\end{equation}
En particular, RH es equivalente a que todos los ceros no triviales pertenezcan a $\operatorname{Fix}(R)$.
\end{corollary}

\subsection{Ejemplo mínimo trabajado: cierre crítico frente a órbita genérica}\label{sec:workedexample}
Para visualizar la clasificación sin depender de una tabla numérica de ceros, considérese primero el punto modelo crítico
\[
\rho_c=\frac12+14i.
\]
Entonces
\[
C(\rho_c)=\frac12-14i,\qquad F(\rho_c)=\frac12-14i,\qquad R(\rho_c)=\rho_c,
\]
y
\[
\Delta_R(\rho_c)=0,
\qquad
\mathcal O(\rho_c)=\left\{\frac12+14i,\frac12-14i\right\}.
\]
El modelo pertenece por construcción al patrón $F_2$: las rutas $C$ y $F$ tienen la misma imagen y $R$ fija el punto.

Como contraste lógico, considérese el punto fuera de línea
\[
\rho_o=0.4+14i.
\]
Se obtiene
\[
C(\rho_o)=0.4-14i,\qquad F(\rho_o)=0.6-14i,\qquad R(\rho_o)=0.6+14i,
\]
\[
\Delta_R(\rho_o)=0.2\neq0,
\]
y la órbita contiene cuatro posiciones distintas,
\[
\mathcal O(\rho_o)=\{0.4\pm14i,\;0.6\pm14i\}.
\]
Este segundo punto es un contraejemplo geométrico al cierre $F_2$, no una afirmación de que sea un cero de $\zeta$. El par de ejemplos muestra qué detecta exactamente el identificador antes de introducir el grafo: igualdad de imágenes en el caso crítico y duplicación transversal en el caso genérico.

\section{Grafo de simetría para identificar las familias}
El primer grafo no usa proximidad numérica; codifica exactamente la acción de las involuciones.

\begin{definition}[Grafo de simetría $\mathcal G_{\rm sym}(\rho)$]
Los vértices son los elementos distintos de $\mathcal O(\rho)$. Se añaden aristas etiquetadas por $C$, $F$ y $R$ entre un vértice $z$ y su imagen bajo cada transformación. Las etiquetas se conservan aun cuando dos transformaciones tengan la misma imagen.
\end{definition}

Para una órbita genérica $F_4$, con
\[
v_1=\rho,\quad v_2=\overline\rho,\quad v_3=1-\rho,\quad v_4=1-\overline\rho,
\]
las aristas son
\begin{align*}
C:&\ (v_1,v_2),(v_3,v_4),\\
F:&\ (v_1,v_3),(v_2,v_4),\\
R:&\ (v_1,v_4),(v_2,v_3).
\end{align*}
Al olvidar etiquetas, el grafo subyacente es $K_4$, con
\begin{equation}
A_{F_4}=\begin{pmatrix}
0&1&1&1\\1&0&1&1\\1&1&0&1\\1&1&1&0
\end{pmatrix},\qquad
L_{F_4}=3I-A_{F_4},
\end{equation}
y espectro
\begin{equation}
 \operatorname{Spec}(L_{F_4})=\{0,4,4,4\}.
\end{equation}

En $F_2$, las identificaciones \eqref{eq:F2collapse} producen dos vértices. Las etiquetas $C$ y $F$ describen la misma transposición, mientras $R$ fija cada vértice. El grafo simple cociente es $K_2$:
\begin{equation}
A_{F_2}=\begin{pmatrix}0&1\\1&0\end{pmatrix},\qquad
L_{F_2}=\begin{pmatrix}1&-1\\-1&1\end{pmatrix},
\end{equation}
con
\begin{equation}
 \operatorname{Spec}(L_{F_2})=\{0,2\}.
\end{equation}
Por tanto, el tamaño de la órbita y el grafo cociente distinguen exactamente el caso crítico del caso genérico. Esta clasificación proporciona el identificador algebraico que se empleará después en el criterio global de celdas.

\begin{figure}[t]
\centering
\begin{tikzpicture}[>=Stealth, every node/.style={circle,draw,minimum size=9mm,font=\small}]
\node (a) at (0,1.5) {$\rho$};
\node (b) at (0,-1.5) {$\bar\rho$};
\node (c) at (3,1.5) {$1-\rho$};
\node (d) at (3,-1.5) {$1-\bar\rho$};
\draw[->] (a) -- node[left,draw=none,fill=white] {$C$} (b);
\draw[->] (a) -- node[above,draw=none,fill=white] {$F$} (c);
\draw[->] (a) -- node[sloped,above,draw=none,fill=white] {$R$} (d);
\draw[->] (c) -- node[right,draw=none,fill=white] {$C$} (d);
\draw[->] (b) -- node[below,draw=none,fill=white] {$F$} (d);
\draw[->] (b) -- node[sloped,below,draw=none,fill=white] {$R$} (c);
\node[draw=none] at (6,0) {$\Longrightarrow\ \Re(\rho)=1/2\ \Longrightarrow$};
\node (e) at (9,1) {$\rho$};
\node (f) at (9,-1) {$\bar\rho$};
\draw[<->,thick] (e) -- node[right,draw=none,fill=white] {$C=F$} (f);
\path[->] (e) edge[loop right] node[right,draw=none] {$R$} (e);
\path[->] (f) edge[loop right] node[right,draw=none] {$R$} (f);
\end{tikzpicture}
\caption{Grafo etiquetado de la órbita genérica $F_4$ y su colapso geométrico a $F_2$ cuando $1-\rho=\bar\rho$. Las etiquetas representan transformaciones distintas, aunque puedan coincidir en su imagen.}
\label{fig:symgraph}
\end{figure}

\section{Operador inducido y conmutación de los dos lenguajes}
Sea $Z^*$ el conjunto de ceros no triviales y sea
\[
Q:Z^*\longrightarrow V(\mathcal G_{\rm sym}),\qquad Q(\rho)=v_{\rho},
\]
la asignación de cada cero a su vértice. La involución $R(s)=1-\overline{s}$ induce una permutación lineal de los vértices,
\begin{equation}\label{eq:PR}
P_R v_{\rho}=v_{R(\rho)}.
\end{equation}
Como $R^2=I$, también
\begin{equation}
P_R^2=I,\qquad P_R^{-1}=P_R.
\end{equation}

\begin{proposition}[Equivarianza del identificador]\label{prop:equivariance}
La asignación $Q$ entre ceros y vértices es equivariante respecto de $R$ y $P_R$:
\begin{equation}\label{eq:commdiagram}
Q\circ R=P_R\circ Q.
\end{equation}
En particular,
\[
Q(\operatorname{Fix}R)=\operatorname{Fix}(P_R)\cap Q(Z^*).
\]
\end{proposition}

\begin{proof}
Por definición, $(Q\circ R)(\rho)=Q(R(\rho))=v_{R(\rho)}$, mientras $(P_R\circ Q)(\rho)=P_Rv_{\rho}=v_{R(\rho)}$. La igualdad de conjuntos fijos es inmediata.
\end{proof}

\begin{proposition}[Conmutación con el Laplaciano]\label{prop:lapcomm}
Si la acción de $P_R$ preserva las aristas del grafo considerado, entonces $P_R$ es un automorfismo y
\begin{equation}\label{eq:lapcomm}
P_RAP_R^{-1}=A,\qquad P_RDP_R^{-1}=D,
\end{equation}
por lo que
\begin{equation}
P_RLP_R^{-1}=L,\qquad [L,P_R]=0,
\end{equation}
con $L=D-A$.
\end{proposition}

\begin{proof}
Un automorfismo preserva adyacencias y grados. Las dos identidades de \eqref{eq:lapcomm} siguen de la representación por matrices de permutación; al restarlas se obtiene la invariancia de $L$. Como $P_R^{-1}=P_R$, la invariancia equivale a $LP_R=P_RL$.
\end{proof}

La conmutación anterior tiene dos lecturas precisas. En $\mathcal G_{\rm sym}$ es una consecuencia exacta de la acción de simetría usada para construir la órbita. En grafos obtenidos por una regla independiente, la igualdad $[L,P_R]=0$ constituye un test estructural de que la arquitectura resultante conserva la misma involución. De esta forma, la simetría compleja y el lenguaje espectral del grafo quedan ligados por el diagrama conmutativo \eqref{eq:commdiagram}.

\section{Generaciones: de las familias primarias a las zetas secundarias}
La palabra \emph{generación} se usa aquí con una separación explícita entre literatura establecida y extensión propuesta.

\subsection{Primera generación: familias de ceros}
La primera generación está formada por las órbitas de los ceros no triviales bajo $C$ y $F$. En la recta crítica, cada familia geométrica es
\begin{equation}
  F_2(\gamma_k)=\left\{\frac12+i\gamma_k,\frac12-i\gamma_k\right\}.
\end{equation}
Al trabajar con $\gamma_k>0$, el cociente por conjugación permite el código escalar
\begin{equation}\label{eq:Q}
 Q(F_2(\gamma_k))=\gamma_k^2.
\end{equation}
Sobre su imagen, \eqref{eq:Q} es biyectivo, con inversa
\begin{equation}
 Q^{-1}(u)=\left\{\frac12+i\sqrt u,\frac12-i\sqrt u\right\},\qquad u>0.
\end{equation}

\subsection{Segunda generación: zetas construidas sobre ceros}
Voros estudió funciones zeta definidas por series de Dirichlet sobre los ceros de Riemann. Una forma fundamental es
\begin{equation}\label{eq:superzeta}
 \mathcal Z(\sigma,v)=\sum_{k=1}^{\infty}(\tau_k^2+v)^{-\sigma},
\end{equation}
con continuación meromorfa en $\sigma$. En la literatura posterior estas construcciones se describen explícitamente como zetas de \emph{segunda generación} o \emph{superzetas}.

La literatura permite separar dos hechos. Voros construye zetas de segunda generación a partir de los ceros, y Arias de Reyna formula una zeta secundaria sin presuponer RH mediante parámetros $\alpha_n$ definidos por $\rho_n=1/2+i\alpha_n$. En el estrato $F_2$, $\alpha_n=\gamma_n\in\mathbb R$ y el código $\gamma_k^2$ es real, unidimensional y reversible sobre la familia crítica; fuera de ese estrato el parámetro puede conservar información transversal compleja. Esta distinción fija con precisión el contenido generacional utilizado en el método.

\subsection{Tercera capa propuesta: familias topológicas persistentes}
Proponemos organizar los códigos positivos $\gamma_k$ o sus coordenadas desplegadas mediante una filtración de proximidad. Esta tercera capa no es una terminología estándar de teoría analítica de números: es una herramienta de clasificación topológica para estudiar si agrupamientos semejantes reaparecen en bloques de ceros a diferentes alturas.

\section{Tres arquitecturas de control: nulo, neutro y no nulo}
Para separar los ceros estructurales que dependen de la arquitectura de los patrones propios de las ordenadas de Riemann se utilizan tres controles con el mismo número $N$ de vértices.

\subsection{Grafo nulo $\mathcal G_{\varnothing}$}
\begin{equation}
 E_{\varnothing}=\varnothing,\qquad A_{\varnothing}=0,\qquad L_{\varnothing}=0.
\end{equation}
Por tanto,
\begin{equation}
 \operatorname{Spec}(L_{\varnothing})=\{0,\ldots,0\},\qquad \beta_0=N.
\end{equation}
Este control representa ausencia completa de conexiones.

\subsection{Grafo neutro $\mathcal G_{\rm neu}(\varepsilon)$}
Se utilizan puntos equiespaciados
\begin{equation}
 y_n=n,\qquad n=1,\ldots,N,
\end{equation}
y la misma regla de proximidad que se usará para los datos:
\begin{equation}\label{eq:neutraledge}
 (i,j)\in E_{\rm neu}(\varepsilon)
 \iff |y_i-y_j|\le\varepsilon.
\end{equation}
Así, el control neutro posee una única escala de separación. Para $\varepsilon=1$ el 1-esqueleto es un camino conectado; la fusión de componentes ocurre esencialmente en una escala única.

\subsection{Grafo no nulo de Riemann $\mathcal G_R(\varepsilon)$}
Cada vértice corresponde a una ordenada positiva $\gamma_n$ de un cero conocido en la recta crítica. Debido al cambio de densidad con la altura, primero se aplica el \emph{unfolding} mediante la parte suave de la fórmula de Riemann--von Mangoldt:
\begin{equation}\label{eq:unfold}
 x_n=\overline N(\gamma_n),\qquad
 \overline N(T)=\frac{T}{2\pi}\log\frac{T}{2\pi}-\frac{T}{2\pi}+\frac78.
\end{equation}
La arista se activa por la misma regla de Rips:
\begin{equation}\label{eq:riemannedge}
 (i,j)\in E_R(\varepsilon)
 \iff |x_i-x_j|\le\varepsilon.
\end{equation}
Definimos
\begin{equation}
 A_R(\varepsilon)=(a_{ij}),\quad
 a_{ij}=\mathbf 1\{0<|x_i-x_j|\le\varepsilon\},
\end{equation}
\begin{equation}
 D_R(\varepsilon)=\operatorname{diag}(d_1,\ldots,d_N),\qquad
 L_R(\varepsilon)=D_R(\varepsilon)-A_R(\varepsilon).
\end{equation}
La multiplicidad del autovalor cero del Laplaciano es el número de componentes conexas:
\begin{equation}\label{eq:beta0}
 \beta_0(\varepsilon)=\dim\ker L_R(\varepsilon).
\end{equation}
Al variar $\varepsilon$, las aristas se activan de forma monótona y se obtiene una filtración de Vietoris--Rips. En dimensión cero, cada muerte de una barra de persistencia corresponde a una fusión de componentes.

\section{Método completo paso a paso}
\begin{algorithmicdef}[Pipeline reproducible]
El procedimiento utilizado para identificar familias y agrupamientos es el siguiente.
\begin{enumerate}[label=\textbf{Paso \arabic*.},leftmargin=*]
\item \textbf{Entrada analítica.} Fijar las transformaciones $C,F,R$ de \eqref{eq:CFR} y la órbita \eqref{eq:orbit}.
\item \textbf{Identificación de familia de simetría.} Construir $\mathcal G_{\rm sym}(\rho)$ y contar posiciones distintas. Si $R(\rho)=\rho$, la familia es $F_2$; si no, la órbita genérica es $F_4$.
\item \textbf{Comprobación algebraica.} Evaluar $\Delta_R=1-2\Re(\rho)$. Cero exacto equivale a la recta crítica.
\item \textbf{Selección de datos.} Tomar un bloque ordenado $\gamma_1<\cdots<\gamma_N$ de ordenadas de ceros verificados. Las tablas de Odlyzko proporcionan, por ejemplo, los primeros $100\,000$ ceros con alta precisión.
\item \textbf{Despliegue.} Transformar $\gamma_n\mapsto x_n$ mediante \eqref{eq:unfold}, para retirar el cambio suave de densidad.
\item \textbf{Control nulo.} Construir $\mathcal G_{\varnothing}$ con los mismos $N$ vértices y sin aristas.
\item \textbf{Control neutro.} Construir $\mathcal G_{\rm neu}(\varepsilon)$ sobre puntos equiespaciados y usar exactamente los mismos valores de $\varepsilon$.
\item \textbf{Grafo no nulo.} Construir $\mathcal G_R(\varepsilon)$ con la regla \eqref{eq:riemannedge}.
\item \textbf{Laplaciano.} Calcular $L=D-A$, su espectro y la multiplicidad de $\lambda=0$ para cada arquitectura.
\item \textbf{Filtración.} Barrer $\varepsilon$ sin escoger un umbral a posteriori. Registrar la curva completa $\beta_0(\varepsilon)$ y los tiempos de muerte $H_0$.
\item \textbf{Familias topológicas.} Definir un agrupamiento a una escala dada como una componente conexa del Rips graph. Las familias más persistentes son aquellas separadas por los mayores tiempos de muerte; no se seleccionan manualmente.
\item \textbf{Comparación de arquitecturas.} Contrastar
\[
\mathcal B_{\varnothing}(\varepsilon),\quad
\mathcal B_{\rm neu}(\varepsilon),\quad
\mathcal B_R(\varepsilon),
\]
donde $\mathcal B=\beta_0$.
\item \textbf{Prueba de autosimilitud estadístico-topológica.} Repetir el pipeline en bloques disjuntos de ceros; comparar las curvas normalizadas $\beta_0(\varepsilon)/N$ y las distribuciones de persistencia. Una coincidencia estable entre bloques sería evidencia de universalidad topológica estadística, no una prueba de RH.
\item \textbf{Capa generacional.} Para familias $F_2$, aplicar el cociente $Q$ y construir, cuando corresponda, objetos secundarios del tipo \eqref{eq:superzeta}. Mantener separada esta generación analítica de la clasificación topológica por proximidad.
\end{enumerate}
\end{algorithmicdef}

\section{Ilustración reproducible con los primeros 300 ceros}
Como control numérico, se calcularon los primeros $300$ ceros positivos mediante aritmética multiprecisión estándar y se verificó la escala contra las tablas públicas de Odlyzko. Después del despliegue \eqref{eq:unfold}, para los $299$ gaps consecutivos se obtiene
\begin{equation}
 \overline d=0.9985241,\qquad
 \operatorname{mediana}(d)=0.9591051,
\end{equation}
\begin{equation}
 d_{\min}=0.2910884,\qquad d_{\max}=1.9934329.
\end{equation}
La Tabla \ref{tab:beta} muestra la curva $H_0$ en ocho escalas predefinidas.

\begin{table}[t]
\centering
\caption{Número de componentes $\beta_0$ para $N=300$ bajo la misma escala $\varepsilon$. El grafo nulo no activa aristas; el neutro equiespaciado fusiona en la escala unitaria; el no nulo de Riemann presenta una transición multiescala.}
\label{tab:beta}
\begin{tabular}{rrrr}
\toprule
$\varepsilon$ & $\beta_0(\mathcal G_{\varnothing})$ & $\beta_0(\mathcal G_{\rm neu})$ & $\beta_0(\mathcal G_R)$\\
\midrule
0.25 & 300 & 300 & 300\\
0.50 & 300 & 300 & 282\\
0.75 & 300 & 300 & 211\\
1.00 & 300 & 1   & 131\\
1.25 & 300 & 1   & 74\\
1.50 & 300 & 1   & 31\\
1.75 & 300 & 1   & 12\\
2.00 & 300 & 1   & 1\\
\bottomrule
\end{tabular}
\end{table}

A $\varepsilon=1.50$ aparecen $31$ componentes; los seis tamaños máximos son
\[
19,\;17,\;16,\;15,\;14,\;12.
\]
A $\varepsilon=1.75$ quedan $12$ componentes; los seis tamaños máximos son
\[
40,\;33,\;32,\;31,\;30,\;26.
\]
Estos grupos son resultados de la filtración $H_0$ y no deben confundirse con las órbitas algebraicas $F_2/F_4$.

En términos espectrales, a $\varepsilon=1$ el grafo de Riemann tiene $169$ aristas y $131$ componentes, mientras el control neutro tiene $299$ aristas y es conectado. A $\varepsilon=2$ el grafo de Riemann es conectado, con $448$ aristas y grado medio $2.9867$. Estos valores funcionan como comprobación de implementación de la arquitectura, no como criterio de verdad de RH.

\section{Ceros del Laplaciano frente a ceros de $\zeta$}
Es esencial mantener separados dos usos de la palabra \emph{cero}.

\begin{enumerate}[label=(\alph*)]
\item Un cero de $\zeta$ o $\xi$ es un punto complejo $\rho$ tal que $\xi(\rho)=0$.
\item Un cero del Laplaciano es un autovalor $\lambda=0$ de $L$; su multiplicidad es $\beta_0$, el número de componentes conexas.
\end{enumerate}

El grafo no ``produce'' ceros de Riemann. Usa ceros previamente calculados como vértices y produce invariantes topológicos de su organización. En cambio, el grafo de simetría $\mathcal G_{\rm sym}$ codifica exactamente cómo cada cero se relaciona con sus imágenes bajo la ecuación funcional y la conjugación.

\section{Polos, inversos y multiplicidad}
La función $\zeta$ tiene un polo simple en $s=1$, mientras $\xi$ es entera. No obstante, los ceros de $\xi$ aparecen como polos de $1/\xi$ y de la derivada logarítmica $\xi'/\xi$.

Si $\rho$ es un cero de multiplicidad $m$, localmente
\begin{equation}
 \xi(s)=(s-\rho)^m g(s),\qquad g(\rho)\neq0,
\end{equation}
y por tanto
\begin{equation}\label{eq:logder}
 \frac{\xi'(s)}{\xi(s)}=\frac{m}{s-\rho}+\frac{g'(s)}{g(s)}.
\end{equation}
Así, $\xi'/\xi$ tiene un polo simple de residuo $m$ en $\rho$.

\begin{proposition}[Principio del argumento y ampliación a simplicidad]
Sea $B$ una región cuya frontera no contiene ceros de $\xi$. Entonces
\begin{equation}
 N(B)=\frac{1}{2\pi i}\oint_{\partial B}\frac{\xi'(s)}{\xi(s)}\,ds
\end{equation}
cuenta ceros con multiplicidad. Si además $R(B)=B$ y $N(B)=1$, el único cero de $B$ está en $\Re(s)=1/2$ y es simple.
\end{proposition}

\begin{proof}
Si $\rho\in B$ es cero, $R(\rho)$ también lo es y pertenece a $B$. Como $N(B)=1$, no puede haber dos posiciones distintas ni multiplicidad mayor que uno; por tanto $R(\rho)=\rho$ y $m=1$. La Proposición \ref{prop:residual} da $\Re(\rho)=1/2$.
\end{proof}

\begin{remark}
La condición $N(B)=1$ contiene dos resultados simultáneos: fija una única posición geométrica dentro de una región $R$-invariante y fija multiplicidad uno. Por ello, una certificación global de celdas unitarias produce una conclusión estrictamente reforzada: localización crítica y simplicidad.
\end{remark}

\section{Mecanismo de comprobación global}\label{sec:globalmechanism}
La ruta lógica que conecta la arquitectura gráfica con el resultado reforzado se resume en
\begin{equation}\label{eq:globalchain}
\boxed{\begin{aligned}
\mathcal G &\Longrightarrow \{B_n\}_{n\ge1}
\Longrightarrow \bigl(R(B_n)=B_n,\;N(B_n)=1\bigr),\\
&\Longrightarrow R(\rho_n)=\rho_n
\Longrightarrow \Re(\rho_n)=\frac12,\;m(\rho_n)=1.
\end{aligned}}
\end{equation}
El contenido de cada paso es el siguiente.
\begin{enumerate}[label=\textbf{G\arabic*.},leftmargin=*]
\item \textbf{Arquitectura.} El grafo de simetría fija las acciones $C,F,R$ y el operador $P_R$; los grafos de proximidad proporcionan una partición multiescala independiente de las ordenadas.
\item \textbf{Celdas.} Una certificación global asocia a cada cero una región de Jordan $B_n$ cuya frontera evita los ceros y cuya arquitectura es estable bajo $R$.
\item \textbf{Índice unitario.} El principio del argumento calcula
\[
N(B_n)=\frac{1}{2\pi i}\oint_{\partial B_n}\frac{\xi'(s)}{\xi(s)}\,ds.
\]
El valor $N(B_n)=1$ afirma simultáneamente una única posición y multiplicidad uno dentro de la celda.
\item \textbf{Fijación por simetría.} Si $R(B_n)=B_n$, la imagen $R(\rho_n)$ permanece en la misma celda; el índice unitario obliga a $R(\rho_n)=\rho_n$.
\item \textbf{Conclusión.} La caracterización crítica da $\Re(\rho_n)=1/2$ y el conteo da $m(\rho_n)=1$.
\end{enumerate}
Así, la condición global que concentra la comprobación es la certificación exhaustiva del segundo paso de \eqref{eq:globalchain}: construir las celdas desde una regla que no introduzca previamente la igualdad $\Re(\rho)=1/2$. El Teorema~\ref{thm:unitcells} formaliza el resultado una vez establecida esa certificación.

\section{Ruta crítica por saturación y cierre por teoría de conjuntos}\label{sec:saturation}
\subsection{Separación entre la lista computacional y el conjunto analítico universal}\label{subsec:list-vs-set}
La demostración universal no se define a partir de una lista previamente tabulada de ordenadas $\gamma_n$. Esa lista pertenece a la capa computacional del método. Para un bloque finito puede escribirse
\begin{equation}\label{eq:computational-list}
 \Gamma_{\rm data}^{(N)}=\{\gamma_1,\ldots,\gamma_N\},
\end{equation}
y utilizarse para desplegar las ordenadas, construir los grafos de proximidad y reproducir los controles numéricos. Su función es comprobar la implementación y estudiar la organización topológica de datos conocidos.

La capa universal parte, en cambio, de la definición analítica del conjunto de ceros. Para una altura $T$ cuya frontera no atraviesa ceros, sea
\begin{equation}\label{eq:analytic-support}
 Z^*(T)=\{\rho:\xi(\rho)=0,\ 0<\Re\rho<1,\ 0<\Im\rho<T\},
\end{equation}
el conjunto de \emph{posiciones} de los ceros no triviales en la región truncada, y sea $m(\rho)\in\mathbb N$ su multiplicidad. El conteo total se define independientemente de cualquier lista finita por
\begin{equation}\label{eq:NTmultiplicity}
 N(T)=\sum_{\rho\in Z^*(T)}m(\rho),
\end{equation}
lo que coincide con el número de ceros contado con multiplicidad obtenido mediante el principio del argumento. Esta separación evita circularidad: los datos tabulados alimentan la comprobación computacional, mientras $Z^*(T)$ y $N(T)$ constituyen los objetos de la afirmación universal.

La lógica de saturación compara entonces un conjunto certificado $S(T)$ con el universo analítico completo $Z^*(T)$, no con una supuesta lista exhaustiva de ceros. En particular, si $S(T)\subseteq Z^*(T)$, todos los elementos de $S(T)$ son simples y
\begin{equation}\label{eq:saturation-simple-chain}
 |S(T)|=N(T)=\sum_{\rho\in Z^*(T)}m(\rho),
\end{equation}
entonces
\[
 |S(T)|\le |Z^*(T)|\le \sum_{\rho\in Z^*(T)}m(\rho)=|S(T)|.
\]
Por consiguiente, ambas desigualdades son igualdades: $S(T)=Z^*(T)$ y $m(\rho)=1$ para todo $\rho\in Z^*(T)$. Esta observación es la forma conjuntista exacta de la saturación y será utilizada en el teorema siguiente.

La formulación por conjuntos reduce el cierre universal a inclusión, conteo y exhaustividad. Sea
Retomando $Z^*(T)$ de \eqref{eq:analytic-support}, definimos el subconjunto fijo
\begin{equation}\label{eq:Z2T}
 Z_2(T)=\{\rho\in Z^*(T):R(\rho)=\rho\}
       =Z^*(T)\cap\{\Re s=1/2\},
\end{equation}
y su complemento relativo
\begin{equation}\label{eq:Z4T}
 Z_4(T)=Z^*(T)\setminus Z_2(T).
\end{equation}
Por construcción,
\begin{equation}\label{eq:partitionT}
 Z^*(T)=Z_2(T)\sqcup Z_4(T).
\end{equation}

Sea $\{B_j(T)\}_{j=1}^{M(T)}$ una colección de celdas certificadas, disjuntas respecto de sus ceros, con
\begin{equation}\label{eq:cellcertT}
 R(B_j)=B_j,\qquad
 N(B_j)=\frac{1}{2\pi i}\oint_{\partial B_j}\frac{\xi'(s)}{\xi(s)}\,ds=1.
\end{equation}
Denotemos por
\begin{equation}\label{eq:ST}
 S(T)=\{\rho_j:\rho_j\in B_j(T),\ j=1,\ldots,M(T)\}
\end{equation}
el conjunto de ceros certificados por dichas celdas. Por la Proposición del principio del argumento,
\begin{equation}\label{eq:inclusionST}
 S(T)\subseteq Z_2(T)\subseteq Z^*(T),
\end{equation}
y cada elemento de $S(T)$ es simple.

\begin{theorem}[Cierre por saturación de conjuntos]\label{thm:saturation}
Supóngase que para una altura $T$ cuya frontera no contiene ceros se tiene la inclusión \eqref{eq:inclusionST} y, además, la condición de saturación
\begin{equation}\label{eq:saturationcount}
 |S(T)|=M(T)=N(T)=\sum_{\rho\in Z^*(T)}m(\rho),
\end{equation}
donde $N(T)$ es el conteo total de ceros no triviales en $0<\Im s<T$, con multiplicidad, definido en \eqref{eq:NTmultiplicity}. Entonces
\begin{equation}\label{eq:saturationequality}
 S(T)=Z_2(T)=Z^*(T),\qquad Z_4(T)=\varnothing,
\end{equation}
y todos los ceros no triviales de altura menor que $T$ están en $\Re(s)=1/2$ y son simples.
\end{theorem}

\begin{proof}
Como $S(T)\subseteq Z^*(T)$, se tiene $|S(T)|\le |Z^*(T)|$. Como toda multiplicidad satisface $m(\rho)\ge1$,
\[
 |Z^*(T)|\le \sum_{\rho\in Z^*(T)}m(\rho)=N(T).
\]
La hipótesis de saturación \eqref{eq:saturationcount} identifica los extremos, de modo que
\[
 |S(T)|=|Z^*(T)|=N(T).
\]
La primera igualdad, junto con $S(T)\subseteq Z^*(T)$ y finitud, da $S(T)=Z^*(T)$. La segunda obliga además a $m(\rho)=1$ para todo $\rho\in Z^*(T)$. Junto con $S(T)\subseteq Z_2(T)\subseteq Z^*(T)$, la doble inclusión produce
\[
 S(T)=Z_2(T)=Z^*(T).
\]
La partición \eqref{eq:partitionT} implica entonces $Z_4(T)=\varnothing$.
\end{proof}

\begin{corollary}[Cierre exhaustivo]\label{cor:exhaustive}
Sea $T_1<T_2<\cdots\to\infty$ una sucesión cuyas fronteras evitan ceros. Si la condición de saturación \eqref{eq:saturationcount} se verifica para cada $T_k$, entonces
\begin{equation}\label{eq:globalsetclosure}
 Z^*=Z_2,\qquad Z_4=\varnothing,
\end{equation}
y, en consecuencia,
\begin{equation}\label{eq:RHsimpleConclusion}
 \Re(\rho)=\frac12,\qquad m(\rho)=1\qquad(\rho\in Z^*).
\end{equation}
\end{corollary}

\begin{proof}
Toda altura finita queda contenida en algún $T_k$. Por el Teorema~\ref{thm:saturation}, en cada región truncada todos los ceros pertenecen a $Z_2(T_k)$ y son simples. Tomando la unión creciente sobre $k$ se obtiene \eqref{eq:globalsetclosure} y \eqref{eq:RHsimpleConclusion}.
\end{proof}

\subsection{Lectura didáctica: ``ruta crítica por saturación''}
La analogía vial resume el conteo sin sustituir la demostración. $N(T)$ representa la capacidad total de posiciones ocupadas por ceros hasta la altura $T$, mientras $M(T)$ cuenta las posiciones ya certificadas por celdas unitarias. Cuando
\[
 M(T)=N(T),
\]
la capacidad queda saturada: no existe una posición adicional disponible para una órbita $F_4$ sin exceder el conteo total. La formulación rigurosa de esta intuición es precisamente el Teorema~\ref{thm:saturation}: inclusión más igualdad de conteo implica igualdad de conjuntos.

\begin{figure}[t]
\centering
\begin{tikzpicture}[>=Stealth,node distance=8mm and 10mm,
 box/.style={draw,rounded corners,align=center,minimum height=9mm,inner sep=4pt,font=\small}]
\node[box] (s) {$S(T)$\\ceros certificados};
\node[box,right=of s] (z2) {$Z_2(T)$\\sector fijo $R$};
\node[box,right=of z2] (zs) {$Z^*(T)$\\todos los ceros};
\draw[->,thick] (s) -- node[above] {$\subseteq$} (z2);
\draw[->,thick] (z2) -- node[above] {$\subseteq$} (zs);
\node[box,below=of z2] (count) {$|S(T)|=N(T)=\sum_{\rho\in Z^*(T)}m(\rho)$};
\draw[->,thick] (count) -- (s);
\draw[->,thick] (count) -- (zs);
\node[box,below=of count] (close) {$S(T)=Z_2(T)=Z^*(T)$\\$Z_4(T)=\varnothing$};
\draw[->,thick] (count) -- (close);
\end{tikzpicture}
\caption{Ruta crítica por saturación. La inclusión de los ceros certificados en el sector fijo, combinada con la saturación del conteo total, fuerza la igualdad de conjuntos y excluye el complemento $Z_4(T)$.}
\label{fig:saturation-route}
\end{figure}

\section{Comprobaciones y contraejemplos}
\subsection{Cero sobre la recta crítica}
Para $\rho=1/2+i\gamma$,
\[
\Delta_R=0,\qquad R(\rho)=\rho,
\]
y la órbita geométrica es $F_2$.

\subsection{Cero hipotético fuera de la recta}
Para
\begin{equation}
 \rho=\frac12+\delta+i\gamma,\qquad\delta\neq0,
\end{equation}
se obtiene
\begin{equation}
 \Delta_R=-2\delta\neq0,
\end{equation}
\begin{equation}
 R(\rho)=\frac12-\delta+i\gamma\neq\rho,
\end{equation}
y la órbita completa es $F_4$. Esta órbita respeta las simetrías conocidas; por ello la simetría, por sí sola, no excluye $F_4$.

\subsection{Cero múltiple sobre $1/2$}
Si un cero en la recta crítica tiene multiplicidad $m>1$, la geometría de la órbita sigue siendo $F_2$ y $R(\rho)=\rho$. La multiplicidad aparece únicamente en \eqref{eq:logder}. Por tanto, cierre geométrico y simplicidad son conceptos distintos.

\subsection{Ceros triviales}
Los ceros triviales $-2n$ pertenecen a $\zeta$ pero no al conjunto de ceros de $\xi$ definido en \eqref{eq:xi}. Además, la reflexión $s\mapsto1-s$ envía $-2n$ a $1+2n$, que no es otro cero de $\zeta$. Por ello no participan en las familias $F_2/F_4$ de ceros no triviales.

\subsection{Control de codificación de $F_4$}
El mapa escalar $Q(F_2)=\gamma^2$ es reversible sobre las familias críticas. Si una familia $F_4$ se parametriza por $(\delta,\gamma)$, proyectarla sólo a $\gamma^2$ pierde $\delta$ y deja de ser inyectivo; un código complejo o bidimensional sí puede conservar esa información. Esta comprobación delimita qué parte de la reversibilidad procede de la familia crítica y qué información adicional debe controlarse en una órbita genérica.

\section{Teoremas de cierre y estatus lógico}
\begin{theorem}[Caracterización crítica]
Para cualquier cero no trivial $\rho$,
\begin{equation}
 \Re(\rho)=\frac12
 \iff R(\rho)=\rho
 \iff F(\rho)=C(\rho)
 \iff \Delta_R(\rho)=0.
\end{equation}
En el semiplano superior, estas condiciones equivalen a que la órbita bajo $R$ tenga una única posición geométrica.
\end{theorem}

\begin{theorem}[Criterio global equivalente a RH]
RH es equivalente a cualquiera de las afirmaciones
\begin{align}
 Z^*&\subseteq\operatorname{Fix}(R),\\
 \Delta_R(\rho)&=0\quad\forall\rho\in Z^*,\\
 |\operatorname{Orb}_R(\rho)|&=1\quad\forall\rho\in Z^*.
\end{align}
\end{theorem}

\begin{theorem}[Criterio reforzado de celdas unitarias]\label{thm:unitcells}
Son equivalentes las siguientes afirmaciones:
\begin{enumerate}[label=(\roman*)]
\item RH es cierta y todos los ceros no triviales son simples.
\item Para cada cero no trivial $\rho_n$ existe una región de Jordan $B_n$ cuya frontera no contiene ceros, tal que $R(B_n)=B_n$, $\rho_n\in B_n$ y
\[
N(B_n)=\frac{1}{2\pi i}\oint_{\partial B_n}\frac{\xi'(s)}{\xi(s)}\,ds=1.
\]
Las regiones pueden elegirse localmente sin otros ceros y la colección cubre el conjunto de ceros no triviales.
\end{enumerate}
\end{theorem}

\begin{proof}
(ii)$\Rightarrow$(i): por el principio del argumento, $N(B_n)=1$ implica un único cero contado con multiplicidad uno. La invariancia $R(B_n)=B_n$ obliga a que la imagen $R(\rho_n)$ permanezca en la misma región; la unicidad da $R(\rho_n)=\rho_n$, y por la Proposición~\ref{prop:residual}, $\Re(\rho_n)=1/2$. Esto vale para toda $n$ y además cada multiplicidad es uno.

(i)$\Rightarrow$(ii): bajo RH, cada $\rho_n$ está en la recta fija de $R$; bajo simplicidad y discreción del conjunto de ceros de la función entera $\xi$, existe un disco suficientemente pequeño centrado en $\rho_n$ que no contiene otros ceros. Al estar el centro en $\Re(s)=1/2$, el disco puede elegirse $R$-invariante. El principio del argumento da entonces $N(B_n)=1$.
\end{proof}

\begin{scope}
El Teorema~\ref{thm:unitcells} convierte el problema universal en un objetivo geométrico-analítico preciso: certificar, mediante la arquitectura construida, la colección exhaustiva de celdas $R$-invariantes con índice unitario. Esta condición es el punto de enlace entre el grafo, el principio del argumento y el resultado reforzado RH más simplicidad.
\end{scope}

\section{Familias arquitectónicas e ``isótopos'' como analogía}
En el análisis exploratorio se utilizó la palabra \emph{isótopo arquitectónico} únicamente como analogía: dentro de una misma naturaleza de grafo se varía un parámetro (por ejemplo, $\varepsilon$, grado medio o tamaño de bloque) mientras se conserva la regla de construcción. Formalmente el objeto es una familia parametrizada
\begin{equation}
 \{\mathcal G_R(\varepsilon)\}_{\varepsilon\ge0}.
\end{equation}
La terminología estándar es \emph{filtración}; ``isótopo'' no se propone como nuevo término técnico.

\section{Reproducibilidad computacional}
La implementación mínima requiere:
\begin{enumerate}
\item una lista ordenada de $\gamma_n$;
\item la función de despliegue \eqref{eq:unfold};
\item matrices de distancia $|x_i-x_j|$;
\item matrices $A(\varepsilon)$, $D(\varepsilon)$ y $L(\varepsilon)$;
\item un barrido de $\varepsilon$ y cálculo de $\beta_0$;
\item opcionalmente, software de homología persistente como GUDHI para verificar el diagrama de persistencia de Vietoris--Rips.
\end{enumerate}
El código suplementario asociado a este manuscrito implementa exactamente estos pasos y reproduce la Tabla \ref{tab:beta}.

\section{Discusión}
El método articula tres niveles complementarios. La prueba universal se formula sobre el conjunto analítico $Z^*(T)$ y el conteo $N(T)$, mientras las listas numéricas de ordenadas se reservan a la capa computacional; esta separación mantiene independiente la exhaustividad lógica de la comprobación numérica. El grafo de simetría es exacto y algebraico: identifica si una órbita concreta pertenece a $F_2$ o $F_4$. El grafo de proximidad es topológico-espectral: cuantifica el agrupamiento de las ordenadas y sus escalas de persistencia. La representación operadora enlaza ambos lenguajes mediante equivarianza y conmutación, mientras la estructura generacional organiza la transmisión de información entre familias primarias y objetos secundarios. La Sección~\ref{sec:correspondencia} amplía esta lectura al comparar el esquema con correspondencias espectrales independientes conocidas en física matemática.

Las zetas secundarias sobre ceros son objetos establecidos en la literatura; la organización de sus códigos mediante persistencia constituye la extensión metodológica aquí desarrollada. El criterio de celdas unitarias especifica además una condición global verificable: cuando la arquitectura certifica $R(B_n)=B_n$ y $N(B_n)=1$ de forma exhaustiva, se obtiene simultáneamente RH y simplicidad. Así, la clasificación orbital, la representación $P_R$, el Laplaciano y el principio del argumento quedan integrados en una misma cadena matemática.


\section{Correspondencia estructural independiente: dos líneas paralelas de lógica}\label{sec:correspondencia}

\subsection{Dos construcciones de origen distinto}
El método permite distinguir dos líneas de construcción que no nacen de la misma pregunta matemática. La primera es analítico--topológica: parte de la función completada de Riemann, sus involuciones y la organización de los ceros,
\begin{equation}\label{eq:lineaA}
 \xi(s)\;\longrightarrow\;R(s)=1-\overline{s}\;\longrightarrow\;
 \operatorname{Fix}(R)\;\longrightarrow\;P_R\;\longrightarrow\;L .
\end{equation}
La segunda es operatorio--espectral: parte de un espacio de estados, una simetría y un generador de dinámica,
\begin{equation}\label{eq:lineaB}
 \mathcal H\;\longrightarrow\;\mathcal J\;\longrightarrow\;
 \operatorname{Fix}(\mathcal J)\;\longrightarrow\;\widehat H\;\longrightarrow\;
 \operatorname{Spec}(\widehat H).
\end{equation}
En la extensión monopolar de la Sección~\ref{sec:monopolo}, por ejemplo, $\mathcal J^2=I$ y $[\widehat H_M,\mathcal J]=0$. En el grafo de simetría, $P_R^2=I$ y, cuando la acción preserva la arquitectura, $[L,P_R]=0$. Por tanto, las dos líneas admiten el mismo esquema lógico abstracto:
\begin{equation}\label{eq:parallelpattern}
 \boxed{\text{involución}\;\to\;\text{sector fijo}\;\to\;\text{conmutación}\;\to\;\text{descomposición espectral}.}
\end{equation}

\begin{figure}[H]
\centering
\begin{tikzpicture}[>=Stealth, node distance=7mm and 5mm, every node/.style={font=\scriptsize,align=center}]
\node[draw,rounded corners] (xi) {$\xi$};
\node[draw,rounded corners,right=of xi] (R) {$R$};
\node[draw,rounded corners,right=of R] (fixR) {$\operatorname{Fix}(R)$};
\node[draw,rounded corners,right=of fixR] (PR) {$P_R$};
\node[draw,rounded corners,right=of PR] (Lap) {$L$};
\draw[->] (xi)--(R); \draw[->] (R)--(fixR); \draw[->] (fixR)--(PR); \draw[->] (PR)--(Lap);

\node[draw,rounded corners,below=12mm of xi] (H) {$\mathcal H$};
\node[draw,rounded corners,right=of H] (J) {$\mathcal J$};
\node[draw,rounded corners,right=of J] (fixJ) {$\operatorname{Fix}(\mathcal J)$};
\node[draw,rounded corners,right=of fixJ] (Ham) {$\widehat H$};
\node[draw,rounded corners,right=of Ham] (Spec) {$\operatorname{Spec}(\widehat H)$};
\draw[->] (H)--(J); \draw[->] (J)--(fixJ); \draw[->] (fixJ)--(Ham); \draw[->] (Ham)--(Spec);

\draw[dashed,<->] (R)--node[right,fill=white] {involución} (J);
\draw[dashed,<->] (fixR)--node[right,fill=white] {sector fijo} (fixJ);
\draw[dashed,<->] (Lap)--node[right,fill=white] {espectro} (Spec);
\end{tikzpicture}
\caption{Dos líneas paralelas de lógica. La fila superior representa la construcción analítico--gráfica asociada a Riemann; la inferior, una construcción operatorio--espectral. Las líneas discontinuas señalan correspondencias de arquitectura, no identidad física entre los objetos.}
\label{fig:parallel-lines}
\end{figure}

\subsection{Antecedente independiente: Riemann, matrices aleatorias y espectros cuánticos}
La aparición de una estructura común entre teoría de números y física espectral tiene un antecedente independiente del presente método. Montgomery estudió la correlación por pares de los ceros de la función zeta \cite{Montgomery1973}; Dyson reconoció que la forma estadística obtenida coincidía con la correlación de autovalores de matrices hermíticas aleatorias, una conexión histórica documentada por el Institute for Advanced Study \cite{IASDyson}. La teoría de matrices aleatorias había sido desarrollada en física nuclear para describir regularidades estadísticas de espectros complejos, y posteriormente la conjetura de Bohigas--Giannoni--Schmit relacionó fluctuaciones espectrales de sistemas cuánticos caóticos con ensembles de matrices aleatorias \cite{BGS1984}. La DLMF recoge de forma explícita las analogías entre la distribución de los ceros de $\zeta$ y autovalores cuánticos semiclasicos, así como el programa Hilbert--P\'olya \cite{DLMFphysical}.

Este antecedente es conceptualmente relevante porque muestra que dos disciplinas construidas desde preguntas distintas pueden converger en leyes espectrales comunes sin que una haya sido definida a partir de la otra. En el lenguaje del presente trabajo, esa convergencia motiva estudiar no sólo coincidencias numéricas, sino morfismos que preserven relaciones.

\subsection{Principio de correspondencia estructural}
Sea $\mathcal A$ una arquitectura analítica y $\mathcal G$ una arquitectura gráfica. Se dirá que existe una \emph{correspondencia estructural} respecto de las involuciones $R$ y $P$ si existe un mapa $Q:\mathcal A\to\mathcal G$ que preserva la acción,
\begin{equation}\label{eq:structcorr}
 \boxed{Q\circ R=P\circ Q.}
\end{equation}
Si, además, $P$ es un automorfismo del grafo, entonces
\begin{equation}\label{eq:structcorrLap}
 P^2=I,\qquad [L,P]=0,
\end{equation}
y el espacio puede descomponerse en sectores propios de $P$. Para la representación de Riemann desarrollada en este manuscrito,
\begin{equation}
 P=P_R,\qquad Q(\operatorname{Fix}R)=\operatorname{Fix}(P_R)\cap Q(Z^*).
\end{equation}
La correspondencia no requiere identificar los objetos físicos con los objetos analíticos. Requiere que las relaciones esenciales sean preservadas. Por ello, el contenido matemático de la comparación reside en la conmutatividad de los diagramas y en la coincidencia de la arquitectura de invariantes, no en una igualdad semántica entre ``cero'', ``partícula'', ``modo'' o ``estado''.

\subsection{Interpretación sobre estructuras no accesibles a observación visual directa}
En física submicroscópica, los grados de libertad elementales y colectivos no se caracterizan por observación visual directa, sino a través de respuestas de detectores, amplitudes de transición, distribuciones de energía y espectros. El modo monopolar proporciona un ejemplo concreto: la coordenada colectiva radial no se observa como una esfera macroscópica que ``respira'', sino que se reconstruye mediante la fuerza de respuesta asociada a un operador como $\widehat M_{\rm IS0}$ y a sus excitaciones espectrales. De forma análoga, los ceros de Riemann se estudian mediante la estructura analítica de $\xi$, sus simetrías, conteos y correlaciones.

La afirmación estructural de este trabajo es, por tanto, la siguiente: arquitecturas que parecen pertenecer a dominios distintos pueden constituir dos líneas paralelas de lógica cuando preservan el mismo patrón de transformación e invariancia. En forma condensada,
\begin{equation}\label{eq:two-logic-lines}
 \boxed{
 \begin{array}{c}
 \text{línea analítica: } R^2=I,\ [L,P_R]=0,\ \operatorname{Fix}(R),\\[1mm]
 \text{línea física: } \mathcal J^2=I,\ [\widehat H,\mathcal J]=0,\ \operatorname{Fix}(\mathcal J).
 \end{array}}
\end{equation}
La coincidencia de forma no se toma como prueba de identidad ontológica; se toma como un objeto matemático investigable. Cuando las dos construcciones son independientes, la persistencia del mismo patrón de involución, punto fijo, conmutación y espectro constituye evidencia de una correspondencia estructural y justifica buscar un formalismo común más general.

\section{Conclusiones}
El marco puede resumirse en cuatro niveles:
\begin{equation}
 \boxed{\text{simetría analítica}\to\text{familia }F_2/F_4\to\text{grafo}\to\text{generación derivada}.}
\end{equation}
El cierre reversible crítico está demostrado por la identidad
\begin{equation}
 1-\rho=\bar\rho\iff\Re(\rho)=1/2.
\end{equation}
El grafo de simetría convierte esta identidad en una reducción de órbita de cuatro a dos posiciones. El pipeline nulo--neutro--no nulo permite estudiar, de forma separada, la topología multiescala de las ordenadas mediante $H_0$. Las funciones secundarias de Voros proporcionan un antecedente establecido para hablar de una segunda generación construida sobre los ceros.

La contribución metodológica del presente texto es integrar estas piezas en un esquema reproducible con una cadena de equivalencias y operadores explícitos. La ruta completa queda sintetizada por las ecuaciones \eqref{eq:roadmap} y \eqref{eq:globalchain}. El Teorema~\ref{thm:saturation} reescribe el cierre universal mediante teoría de conjuntos: $S(T)\subseteq Z_2(T)\subseteq Z^*(T)$ y la saturación $|S(T)|=N(T)$ fuerzan $S(T)=Z_2(T)=Z^*(T)$ y, por tanto, $Z_4(T)=\varnothing$. El Corolario~\ref{cor:exhaustive} extiende este cierre a $T\to\infty$. El criterio reforzado de celdas unitarias muestra que la certificación global de $R(B_n)=B_n$ y $N(B_n)=1$ no sólo sitúa todos los ceros no triviales sobre la recta crítica, sino que simultáneamente establece su simplicidad. La representación inducida $P_R$ y su conmutación con el Laplaciano proporcionan el lenguaje operador que conecta la simetría compleja con la arquitectura del grafo. La correspondencia estructural independiente formula además un segundo nivel de lectura: el mismo patrón abstracto de involución, sector fijo, conmutación y espectro aparece en construcciones físicas operatorias desarrolladas desde problemas distintos. El antecedente Montgomery--Dyson y las conexiones con matrices aleatorias muestran que este tipo de convergencia entre teoría de números y física espectral posee una historia matemática propia; el presente método la expresa mediante diagramas conmutativos e invariantes explícitos.

\appendix
\section{Pseudocódigo del grafo no nulo}
\begin{verbatim}
input: gamma[1..N], epsilon_grid
x[n] = Nbar(gamma[n])

for epsilon in epsilon_grid:
    A = zeros(N,N)
    for i < j:
        if abs(x[i]-x[j]) <= epsilon:
            A[i,j] = A[j,i] = 1
    D = diag(row_sums(A))
    L = D - A
    eigenvalues = eigvalsh(L)
    beta0 = count(abs(eigenvalues) < tolerance)
    save(epsilon, beta0, eigenvalues)
\end{verbatim}

\section{Pseudocódigo del identificador de familias}
\begin{verbatim}
input: rho
C = conjugate(rho)
F = 1 - rho
R = 1 - conjugate(rho)
Delta = F - C

orbit = unique({rho, C, F, R})
if R == rho:
    family = F2
else:
    family = F4
\end{verbatim}

\section{Extensión operatoria: monopolo radial, modo de respiración y espectro}\label{sec:monopolo}
\subsection{Dos significados de ``monopolo''}
En este trabajo el término \emph{operador monopolar} se utiliza en el sentido nuclear de multipolaridad $L=0$. No debe confundirse con el monopolo magnético de Dirac, que corresponde a una carga magnética aislada hipotética. La distinción es relevante: el primero es un operador colectivo establecido en física nuclear; el segundo es una partícula hipotética buscada experimentalmente.

\subsection{Operador monopolar isoescalar}
Para un núcleo de número másico $A$, el operador monopolar isoescalar de una partícula se representa por la coordenada radial cuadrática y, a nivel de muchos cuerpos, por
\begin{equation}\label{eq:mis0}
  \widehat M_{\rm IS0}=\sum_{i=1}^{A} r_i^2 .
\end{equation}
Es un operador escalar bajo rotaciones y selecciona excitaciones monopolares $J^\pi=0^+$. La función de fuerza asociada puede escribirse como
\begin{equation}\label{eq:monostrength}
 S_{\rm IS0}(E)=\sum_{\nu}
 \left|\langle \Psi_\nu|\widehat M_{\rm IS0}|\Psi_0\rangle\right|^2
 \delta\!\left(E-(E_\nu-E_0)\right).
\end{equation}
La resonancia monopolar gigante isoescalar (ISGMR) corresponde a una respuesta colectiva radial en la que protones y neutrones oscilan aproximadamente en fase. Por ello se denomina \emph{breathing mode}: la distribución nuclear se desplaza entre configuraciones de mayor y menor densidad conservando, en primera aproximación, la simetría esférica.

Para transiciones monopolares eléctricas $E0$, el operador de carga contiene el análogo protónico
\begin{equation}\label{eq:e0op}
  \widehat M(E0)=e\sum_{p=1}^{Z} r_p^2,
\end{equation}
y la intensidad de transición está determinada por su elemento de matriz. En una transición $E0$ pura, la emisión de un solo fotón está prohibida por conservación del momento angular; los canales incluyen conversión interna y, cuando la energía lo permite, creación interna de pares $e^-e^+$.

\subsection{Coordenada colectiva y máximo--mínimo}
La geometría del modo puede idealizarse mediante una coordenada colectiva radial
\begin{equation}\label{eq:radialcollective}
  R_c(t)=R_0+A_M\cos(\omega_M t+\phi),
\end{equation}
con
\begin{equation}
  R_{\max}=R_0+A_M,
  \qquad
  R_{\min}=R_0-A_M.
\end{equation}
La variable centrada
\begin{equation}
  x(t)=R_c(t)-R_0
\end{equation}
posee la involución radial
\begin{equation}\label{eq:radialinvolution}
  \mathcal J:x\mapsto -x,
  \qquad
  \mathcal J^2=I.
\end{equation}
Así, expansión y compresión aparecen como estados homólogos respecto del punto fijo $x=0$, que representa el radio de equilibrio. Los extremos $+A_M$ y $-A_M$ forman una pareja bajo $\mathcal J$.

\subsection{Hamiltoniano colectivo y cuantización}
Cerca del equilibrio, la aproximación armónica conduce al Hamiltoniano colectivo
\begin{equation}\label{eq:monoham}
 H_M=\frac{P_c^2}{2M_{\rm eff}}
     +\frac12 C_M(R_c-R_0)^2,
 \qquad
 \omega_M=\sqrt{\frac{C_M}{M_{\rm eff}}},
\end{equation}
donde $M_{\rm eff}$ es la inercia colectiva y $C_M$ la rigidez restauradora. Las ecuaciones de Hamilton implican
\begin{equation}
  \ddot R_c+\omega_M^2(R_c-R_0)=0,
\end{equation}
que recupera \eqref{eq:radialcollective}. Al cuantizar la coordenada radial,
\begin{equation}\label{eq:monohamq}
 \widehat H_M=
 -\frac{\hbar^2}{2M_{\rm eff}}\frac{d^2}{dR_c^2}
 +\frac12 C_M(R_c-R_0)^2,
\end{equation}
y, en la aproximación armónica ideal,
\begin{equation}
 E_n=\hbar\omega_M\left(n+\frac12\right).
\end{equation}
La frecuencia de resonancia determina la escala a la que la respuesta monopolar se maximiza; el ancho espectral observado codifica el amortiguamiento y la fragmentación del modo colectivo.

\subsection{Correspondencia estructural con el lenguaje de operadores del grafo}
El operador monopolar nuclear $\widehat M_{\rm IS0}$ y el operador de simetría $P_R$ pertenecen a teorías distintas, pero permiten comparar una misma arquitectura abstracta:
\begin{equation}\label{eq:operatorparallel}
\begin{array}{ccc}
\text{modo radial} &:& \mathcal J^2=I,\quad x\leftrightarrow -x,\quad \operatorname{Fix}(\mathcal J)=\{0\},\\[1mm]
\text{grafo de Riemann} &:& P_R^2=I,\quad v_\rho\leftrightarrow v_{R(\rho)},\quad \operatorname{Fix}(P_R).
\end{array}
\end{equation}
En ambos casos existe una involución, un conjunto fijo y una descomposición del espacio en sectores pares e impares respecto del operador. Si el operador conmuta con el generador espectral correspondiente, sus autoespacios pueden organizarse simultáneamente. Para el grafo, esta condición es
\begin{equation}
 [L,P_R]=0.
\end{equation}
Para el oscilador monopolar ideal, la reflexión radial $\mathcal J$ conmuta con el Hamiltoniano armónico,
\begin{equation}\label{eq:monoparity}
 [\widehat H_M,\mathcal J]=0,
\end{equation}
porque el potencial depende de $(R_c-R_0)^2$. Por tanto, ambos lenguajes presentan de forma explícita una simetría involutiva compatible con un operador espectral.

Esta correspondencia no identifica físicamente $P_R$ con $\widehat M_{\rm IS0}$ ni con $\widehat H_M$; establece una equivalencia de arquitectura operatoria: involución, punto fijo, pares homólogos, conmutación con un operador espectral y organización de modos propios. Tal paralelismo proporciona un lenguaje común para describir el cierre máximo--mínimo del modo monopolar y el cierre de las órbitas bajo $R(s)=1-\overline{s}$.

\section{Formato editorial}
El manuscrito usa la clase \texttt{amsart}, teoremas numerados, resumen, palabras clave, clasificación MSC y bibliografía convencional, siguiendo la estructura recomendada por el \emph{AMS Author Handbook}. El formato de una revista concreta puede exigir macros, estilo bibliográfico o metadatos adicionales; esos cambios son editoriales y no afectan la matemática.

\begin{thebibliography}{99}
\bibitem{Ahlfors}
L. V. Ahlfors, \emph{Complex Analysis}, 3rd ed., McGraw--Hill, 1979.

\bibitem{Arias}
J. Arias de Reyna, \emph{Computation of the secondary zeta function}, arXiv:2006.04869, 2020.

\bibitem{Bombieri}
E. Bombieri, \emph{The Riemann Hypothesis}, in: The Millennium Prize Problems, Clay Mathematics Institute/AMS, 2006.

\bibitem{ClayRH2026}
Clay Mathematics Institute, \emph{Riemann Hypothesis}, Millennium Prize Problem overview and official resources, accessed 26 August 2026.

\bibitem{DLMFreflection}
NIST Digital Library of Mathematical Functions, \S25.4, \emph{Reflection Formulas: Riemann Zeta Function}, accessed 2026.

\bibitem{DLMFzeros}
NIST Digital Library of Mathematical Functions, \S25.10, \emph{Zeros: Riemann Zeta Function}, accessed 2026.

\bibitem{EdelsbrunnerHarer}
H. Edelsbrunner and J. Harer, \emph{Computational Topology: An Introduction}, American Mathematical Society, 2010.

\bibitem{GUDHI}
GUDHI Project, \emph{Rips complex and persistent homology documentation}, INRIA, accessed 2026.

\bibitem{GodsilRoyle}
C. Godsil and G. Royle, \emph{Algebraic Graph Theory}, Graduate Texts in Mathematics 207, Springer, 2001.

\bibitem{DuguetMonopole}
T. Duguet et al., \emph{Ab initio description of monopole resonances in light- and medium-mass nuclei}, Eur. Phys. J. A (2024), monopole operator $r^2=\sum_i r_i^2$ and radial response.

\bibitem{Harakeh}
M. N. Harakeh and A. van der Woude, \emph{Giant Resonances: Fundamental High-Frequency Modes of Nuclear Excitation}, Oxford University Press, 2001.

\bibitem{IAEAE0}
T. Kib\'edi et al., \emph{BrIcc: Conversion coefficients and electric monopole transitions}, Nuclear Data Section, International Atomic Energy Agency; E0 transition strength and internal-conversion/pair channels.

\bibitem{CERNMonopole}
CERN, \emph{MoEDAL--MAPP and the search for magnetic monopoles}, CERN experimental overview, accessed 2026.

\bibitem{Odlyzko}
A. M. Odlyzko, \emph{Tables of zeros of the Riemann zeta function}, University of Minnesota, electronic tables, accessed 2026.

\bibitem{Riemann}
B. Riemann, \emph{Ueber die Anzahl der Primzahlen unter einer gegebenen Grösse}, Monatsberichte der Berliner Akademie, 1859.

\bibitem{Titchmarsh}
E. C. Titchmarsh, revised by D. R. Heath-Brown, \emph{The Theory of the Riemann Zeta-Function}, 2nd ed., Oxford University Press, 1986.

\bibitem{Voros2003}
A. Voros, \emph{Zeta functions for the Riemann zeros}, Ann. Inst. Fourier 53 (2003), no. 3, 665--699, doi:10.5802/aif.1955.

\bibitem{Voros2010}
A. Voros, \emph{Zeta Functions over Zeros of Zeta Functions}, Lecture Notes of the Unione Matematica Italiana 8, Springer, 2010.


\bibitem{Montgomery1973}
H. L. Montgomery, \emph{The pair correlation of zeros of the zeta function}, in: Analytic Number Theory (St. Louis, 1972), Proc. Sympos. Pure Math. 24, American Mathematical Society, Providence, 1973, pp. 181--193, doi:10.1090/pspum/024/9944.

\bibitem{BGS1984}
O. Bohigas, M. J. Giannoni and C. Schmit, \emph{Characterization of chaotic quantum spectra and universality of level fluctuation laws}, Phys. Rev. Lett. 52 (1984), 1--4, doi:10.1103/PhysRevLett.52.1.

\bibitem{DLMFphysical}
NIST Digital Library of Mathematical Functions, \S25.17, \emph{Physical Applications: Zeta and Related Functions}, accessed 2026.

\bibitem{IASDyson}
Institute for Advanced Study, \emph{From Prime Numbers to Nuclear Physics and Beyond}, historical account of the Montgomery--Dyson correspondence, 2013, accessed 2026.

\bibitem{AMSHandbook}
American Mathematical Society, \emph{AMS Author Handbook for Journals}, current online edition, accessed 2026.
\end{thebibliography}

\end{document}
